\input{../tex-inputs/latex-leseansicht-vorspann}

\section[Gustav Schwarzkopf an Arthur Schnitzler, 10. 8. 1895]{L04280 Gustav Schwarzkopf an Arthur Schnitzler, 10. 8. 1895}
\nopagebreak\mylabel{L04280v}
\rehead{ }\normalsize\beginnumbering\briefempfaengerindex{Schnitzler, Arthur@\textsc{Schnitzler, Arthur}!zzzSchwarzkopf, Gustav@\emph{von Gustav Schwarzkopf}!1895-08-102@{10. 8. 1895}|(be}
\toendnotes[C]{\smallbreak\pagebreak[2]}
\correspDesc{Versand  durch Gustav Schwarzkopf am 10. 8. 1895 in Wien
\newline{}Erhalt  durch Arthur Schnitzler im Zeitraum [11. 8. 1895 – 15. 8. 1895?] in Bad Ischl}\toendnotes[C]{\smallbreak}
\Standort{CUL, Schnitzler, B 96a.}
\physDesc{Briefkarte, 1005 Zeichen
\newline{}Handschrift: schwarze Tinte, deutsche Kurrent}\toendnotes[C]{\smallbreak}
\pstart
           \raggedleft{}{\pb}10. VIII 95\pend
           \vspace{0.5em}
\pstart
           Bitte, lieber Freund, entſchuldigen Sie meine Unhöflichkeit, die Sie
               länger als 14 Tage auf Anwort warten ließ. Mit Überbürdung ka{\geminationn} ich dieſe Verzögerung freilich nicht motiviren; ich
               habe nur i{\geminationm}er auf etwas Mitteilenswertes gewartet, auf
               ein intereſſantes Abenteuer, oder wenigſtens auf ein »inneres Erlebniß«, aber es hat{ }ſich nichts dergleichen eingeſtellt. Wie Sie{ }ſehen, bin {\pb}ich noch i{\geminationm}er in Wien\oindex{Wien@\textbf{Wien}, \emph{Verwaltungsgebiet}|pw}. Bei \label{K_L04280-1v}\edtext{\textsc{Weiss\pwindex{Karlweis, Carl 23.\,11.\,1850 Wien – 27.\,10.\,1901 ebd.@\textsc{Karlweis, Carl} (23.\,11.\,1850 Wien – 27.\,10.\,1901 ebd.), \emph{Schriftsteller}|pw}} iſt zwar eine Freikarte}{\lemma{\textnormal{\emph{Weiss … Freikarte}}}\Cendnote{\textnormal{C. Karlweis\pwindex{Karlweis, Carl 23.\,11.\,1850 Wien – 27.\,10.\,1901 ebd.@\textsc{Karlweis, Carl} (23.\,11.\,1850 Wien – 27.\,10.\,1901 ebd.), \emph{Schriftsteller}|pwk} arbeitete
                  im Zivilberuf für die \emph{Südbahn-Gesellschaft}\orgindex{Südbahn-Gesellschaft@Südbahn-Gesellschaft|pwk} und konnte hier
               anscheinend Fahrkarten bereitstellen.}}}\label{K_L04280-1} für mich vorgemerkt, aber
               ich habe{ }ſie noch nicht abgeholt, werde es auch wahrſcheinlich nicht. Ich hab{ }ſo alle
               Luſt verloren und beruhige mich mit dem Bewußtſein, daß ich bis \textsc{Villach\oindex{Villach@\textbf{Villach}, \emph{Verwaltungsgebiet}|pw}} fahren kö{\geminationn}te, we{\geminationn} ich
               wollte. Das iſt auch etwas. – Von \textsc{R. Golm\pwindex{Goldscheid, Rudolf 12.\,8.\,1870 Wien – 31.\,10.\,1931 ebd.@\textsc{Goldscheid, Rudolf} (12.\,8.\,1870 Wien – 31.\,10.\,1931 ebd.), \emph{Philosoph, Publizist}|pw}} hab ich nichts geleſen. Übrigens hat der Ma{\geminationn} mit{ }ſeinem Benehmen Recht, we{\geminationn} er dadurch Exemplare abſetzt.
                  We{\geminationn} er{ }ſich Ihrem \textsc{Cousin\pwindex{Mandl, Alfred 28.\,2.\,1878 Wien – 11.\,2.\,1926 Prag@\textsc{Mandl, Alfred} (28.\,2.\,1878 Wien – 11.\,2.\,1926 Prag), \emph{Ingenieur}|pwv}}{ }ſympathiſch gemacht hätte, wäre ihm das vielleicht nie gelungen. Ärger zu
               erregen erſcheint alſo ein gutes Mittel zu{ }ſein. Das ka{\geminationn}
               man{ }ſich ja merken.\pend
           \pstart Herzlichſt Ihr \spacefill\mbox{GustavSchwk}\pend{}
\pstart
           \noindent{}Grüße an \textsc{Beer-Hofma{\geminationn}\pwindex{Beer-Hofmann, Richard 11.\,7.\,1866 Wien – 26.\,9.\,1945 New York City@\textsc{Beer-Hofmann, Richard} (11.\,7.\,1866 Wien – 26.\,9.\,1945 New York City), \emph{Schriftsteller}|pw}}.\pend
           \selectlanguage{ngerman}\endnumbering\briefempfaengerindex{Schnitzler, Arthur@\textsc{Schnitzler, Arthur}!zzzSchwarzkopf, Gustav@\emph{von Gustav Schwarzkopf}!1895-08-102@{10. 8. 1895}|)be}\mylabel{L04280h}
\begin{anhang}
\end{anhang}\newcommand{\dateiname}{L04280}\newcommand{\titel}{Gustav Schwarzkopf an Arthur Schnitzler, 10. 8. 1895}\newcommand{\editorInnen}{Herausgegeben von Jahnke, SelmaMüller, Martin Anton}\input{../tex-inputs/latex-leseansicht-abspann}
